% TOP_PROBLEM: {"id": 3065, "title": "Aharoni-Berger conjecture", "category_id": 2, "category": {"id": 2, "name": "combinatorics", "display_name": "Combinatorics", "description": "Counting problems, graph theory, discrete structures.", "slug": "combinatorics", "order_index": 2, "created_at": "2026-07-31T15:26:25.671Z"}, "status": "open", "problem_number": "OPG-382", "set_id": 12, "statusQualification": "Scope completed with k at least two, positive target size, finite ground set, and ordered partitions allowing empty parts."}
% FIRST_POSTED: 2026-10-01
% ATTEMPT_STATUS: unresolved
% UnsolvedMath ID 3065; source catalogue code OPG-382
% !TeX program = lualatex
% GPT-6 Astra Ultra research attempt; independent partial work, 2026-10-01.
\documentclass[11pt,a4paper]{article}
\usepackage[margin=25mm]{geometry}
\usepackage{fontspec}
\setmainfont{lmroman10-regular.otf}[BoldFont=lmroman10-bold.otf,
 ItalicFont=lmroman10-italic.otf,BoldItalicFont=lmroman10-bolditalic.otf]
\usepackage{amsmath,amssymb,mathtools}
\usepackage{unicode-math}
\setmathfont{latinmodern-math.otf}
\AtBeginDocument{\let\setminus\smallsetminus}
\usepackage{xurl,hyperref}
\hypersetup{colorlinks=true,urlcolor=blue}
\setcounter{secnumdepth}{0}
\setlength{\parindent}{0pt}
\setlength{\parskip}{5pt}
\setlength{\emergencystretch}{3em}
\begin{document}
\section*{Aharoni-Berger conjecture}\label{3065}
{\small UnsolvedMath ID 3065 · OPG-382 · Combinatorics · OpenGarden}
\subsection{Definitions and mathematical statement}
Let $E$ be finite and let $k\ge2$ and $\ell\ge1$ be integers.
For each $i\in\{1,\ldots,k\}$, let $M_i=(E,\mathcal I_i)$ be a matroid:
$\mathcal I_i$ is a nonempty hereditary family with the exchange property
that if $I,J\in\mathcal I_i$ and $|I|<|J|$, an element of $J\setminus I$
can be added to $I$ while preserving membership in $\mathcal I_i$.
Its rank function is
$r_i(X)=\max\{|I|:I\subseteq X,\ I\in\mathcal I_i\}$.

An ordered partition of $E$ into $k$ parts means pairwise disjoint sets
$X_1,\ldots,X_k$ with union $E$; empty parts are permitted, and the part
$X_i$ is assigned to the particular matroid $M_i$.
Suppose that every such partition satisfies
\[
 \sum_{i=1}^k r_i(X_i)\ge\ell(k-1).
\]
The conjecture concludes that there exists a single subset $I\subseteq E$
with $|I|=\ell$ and $I\in\bigcap_{i=1}^k\mathcal I_i$.
Thus one selects exactly the same elements independently in all the
matroids, rather than separate independent subsets for the different
parts of a partition.

The factor $k-1$ and the quantification over all ordered partitions are
part of the proposed sufficient condition. The case $k=1$ is excluded:
with that value the displayed hypothesis would be vacuous and could
not imply the claimed conclusion. Allowing empty parts prevents a
vacuous hypothesis when the ground set is small.
\subsection{Short English statement}
Does the specified rank bound over every partition force an independent set of size $\ell$ common to all $k$ matroids?
\subsection{Research attempt: a closure certificate and the missing factor}
\paragraph{Scope and literature check.}
GPT-6 Astra Ultra, 1 October 2026. The parts of the partition are
ordered and can be empty. The original discussion [S2] identifies
$k=2$ with matroid intersection and records the $k=3$ result of
Aharoni--Berger. The authors' research report [S3] explains the
matroid/complex intersection framework. Our elementary argument below
does not reproduce that topological theorem. It obtains a $k$-factor
bound and, by integrality, settles the stated conclusion when $\ell<k$.

\paragraph{1. Define the two extremal quantities.}
Let
\[
 R=\min_{E=X_1\sqcup\cdots\sqcup X_k}\sum_i r_i(X_i),\qquad
 \nu=\max\{|I|: I\text{ is independent in every }M_i\}.
\]
These extrema exist on the finite ground set. If $I$ is common
independent, then $I\cap X_i$ is independent in $M_i$ and hence
$r_i(X_i)\ge|I\cap X_i|$. Summing gives $R\ge\nu$. This
necessary inequality alone has the wrong direction to force a large $I$.

\paragraph{2. Every maximal common independent set certifies $R\le k|I|$.}
Construct $I$ greedily until no element can be added while retaining
independence in all matroids. Write $s=|I|$; maximality here is under
inclusion, not necessarily maximum cardinality. For each $e\notin I$
choose an index $i(e)$ such that $I\cup\{e\}$ is dependent in $M_i$.
Since $I$ is independent, this says
$r_i(I\cup\{e\})=r_i(I)=s$, or $e\in\operatorname{cl}_i(I)$.

For completeness, adjoining all such individually spanned elements
does not raise rank. If $r_i(I\cup A)=s$ and
$r_i(I\cup\{e\})=s$, submodularity yields
\[
 r_i(I\cup A\cup\{e\})+r_i(I)
 \le r_i(I\cup A)+r_i(I\cup\{e\})=2s.
\]
Induction establishes the claim for every finite collection of them.

Assign each outside element to $X_{i(e)}$, and assign the elements of
$I$ arbitrarily to the parts, for example all to $X_1$. Then
$X_i\subseteq\operatorname{cl}_i(I)$, so $r_i(X_i)\le s$ for
every $i$. This is an actual ordered partition, and therefore
\[
 R\le\sum_i r_i(X_i)\le ks.
\]
Applying the hypothesis $R\ge\ell(k-1)$ gives
\[
 |I|\ge\left\lceil\frac{\ell(k-1)}k\right\rceil.
\]
In particular if $1\le\ell<k$, the number inside the ceiling is
$\ell-\ell/k\in(\ell-1,\ell)$, so $|I|\ge\ell$. Any
$\ell$-subset of $I$ answers the question in this parameter range.

\paragraph{3. Where the tempting improvement breaks.}
One might try to replace $ks$ by $(k-1)s$ by saying that each element
of $I$ was assigned to only one part. That does not bound the ranks
of the other parts: a part disjoint from $I$ may still have rank $s$
while lying in its closure. For a rank-one matroid with two parallel
nonloops $a,b$, the set $\{b\}$ is disjoint from the basis $\{a\}$,
lies in its closure, and has full rank one. In higher rank, parallel
copies of all basis elements give the same obstruction. Thus accounting
for membership in $I$ cannot save one whole rank contribution.

\paragraph{Verification and remaining gap.}
The argument permits loops; they simply belong to the closure of every
set and contribute no rank. If $I$ were empty, the constructed partition
would have total rank zero, contradicting the positive hypothesis.
At $\ell=k$, the bound supplies only $k-1$, so the integrality argument
has a genuine endpoint. A proof for larger $\ell$ needs a more effective
exchange or partition certificate than single-element maximality. No
unjustified $k$-matroid analogue of the two-matroid min--max theorem is used.

\subsection{Final status}
UNRESOLVED in general. A constructive lower bound for every maximal
common independent set and the full implication for $\ell<k$ are
proved. The known $k=2,3$ results are attributed to the literature;
no new general solution or completion percentage is claimed.

\subsection{Sources}
\begin{itemize}
\item[S1] ULAM.AI and the UnsolvedMath Contributors, supplied problems.json, record 3065 (OPG-382), OpenGarden collection. Catalogue identity and source transcription; CC BY 4.0.
\url{https://huggingface.co/datasets/ulamai/UnsolvedMath}.
\item[S2] Open Problem Garden, Aharoni-Berger conjecture. Original problem and discussion; the mathematical formulation below is expanded from the supplied source record.
\url{http://www.openproblemgarden.org/op/aharoni_berger_conjecture}.
\item[S3] R. Aharoni, joint work with E. Berger, \emph{The intersection of a matroid and a simplicial complex}, research report in Oberwolfach Report 3/2005, Graph Theory. Author's report checked 1 October 2026.
\url{https://ems.press/content/serial-article-files/45977?nt=1}.
\end{itemize}
\end{document}
